% Exam Template for UMTYMP and Math Department courses
%
% Using Philip Hirschhorn's exam.cls: http://www-math.mit.edu/~psh/#ExamCls
%
% run pdflatex on a finished exam at least three times to do the grading table on front page.
%
%%%%%%%%%%%%%%%%%%%%%%%%%%%%%%%%%%%%%%%%%%%%%%%%%%%%%%%%%%%%%%%%%%%%%%%%%%%%%%%%%%%%%%%%%%%%%%

% These lines can probably stay unchanged, although you can remove the last
% two packages if you're not making pictures with tikz.
\documentclass[11pt]{exam}
\RequirePackage{amssymb, amsfonts, amsmath, latexsym, verbatim, xspace, setspace}
\RequirePackage{tikz, pgflibraryplotmarks}

% By default LaTeX uses large margins.  This doesn't work well on exams; problems
% end up in the "middle" of the page, reducing the amount of space for students
% to work on them.
\usepackage[margin=1in]{geometry}
\usepackage{enumerate}
\usepackage{amsthm}

\theoremstyle{definition}
\newtheorem{soln}{Solution}

% Here's where you edit the Class, Exam, Date, etc.
\newcommand{\class}{Math 407 Section 1}
\newcommand{\term}{Fall 2025}
\newcommand{\examnum}{Exam 2}
\newcommand{\examdate}{March 12, 2025}
\newcommand{\timelimit}{1 Hour 50 Minutes}
\newcommand{\ol}[1]{\overline{#1}}

% For an exam, single spacing is most appropriate
\singlespacing
% \onehalfspacing
% \doublespacing

% For an exam, we generally want to turn off paragraph indentation
\parindent 0ex

\begin{document} 

% These commands set up the running header on the top of the exam pages
\pagestyle{head}
\firstpageheader{}{}{}
\runningheader{\class}{\examnum\ - Page \thepage\ of \numpages}{\examdate}
\runningheadrule

\begin{flushright}
\begin{tabular}{p{2.8in} r l}
\textbf{\class} & \textbf{Name (Print):} & \makebox[2in]{\hrulefill}\\
\textbf{\term} &&\\
\textbf{\examnum} & \textbf{Student ID:}&\makebox[2in]{\hrulefill}\\
\textbf{\examdate} &&\\
\textbf{Time Limit: \timelimit} % & Teaching Assistant & \makebox[2in]{\hrulefill}
\end{tabular}\\
\end{flushright}
\rule[1ex]{\textwidth}{.1pt}


This exam contains \numpages\ pages (including this cover page) and
\numquestions\ problems.  Check to see if any pages are missing.  Enter
all requested information on the top of this page, and put your initials
on the top of every page, in case the pages become separated.\\

You may \textit{not} use your books or notes on this exam.\\

You are required to show your work on each problem on this exam.  The following rules apply:\\

\begin{minipage}[t]{3.7in}
\vspace{0pt}
\begin{itemize}

%\item \textbf{If you use a ``fundamental theorem'' you must indicate this} and explain
%why the theorem may be applied.

\item \textbf{Organize your work}, in a reasonably neat and coherent way, in
the space provided. Work scattered all over the page without a clear ordering will 
receive very little credit.  

\item \textbf{Mysterious or unsupported answers will not receive full
credit}.  A correct answer, unsupported by calculations, explanation,
or algebraic work will receive no credit; an incorrect answer supported
by substantially correct calculations and explanations might still receive
partial credit.  This especially applies to limit calculations.

\item If you need more space, use the back of the pages; clearly indicate when you have done this.

\item If the problem asks for a proof, be sure to carefully justify your work, including any theorems from class.

Note: you may NOT use a theorem or result from class to prove something when it makes the problem entirely trivial.  If you are unsure whether a particular theorem or result is allowed, just ask!

\end{itemize}

Do not write in the table to the right.
\end{minipage}
\hfill
\begin{minipage}[t]{2.3in}
\vspace{0pt}
%\cellwidth{3em}
\gradetablestretch{2}
\vqword{Problem}
\addpoints % required here by exam.cls, even though questions haven't started yet.	
\gradetable[v]%[pages]  % Use [pages] to have grading table by page instead of question

\end{minipage}
\newpage % End of cover page

%%%%%%%%%%%%%%%%%%%%%%%%%%%%%%%%%%%%%%%%%%%%%%%%%%%%%%%%%%%%%%%%%%%%%%%%%%%%%%%%%%%%%
%
% See http://www-math.mit.edu/~psh/#ExamCls for full documentation, but the questions
% below give an idea of how to write questions [with parts] and have the points
% tracked automatically on the cover page.
%
%
%%%%%%%%%%%%%%%%%%%%%%%%%%%%%%%%%%%%%%%%%%%%%%%%%%%%%%%%%%%%%%%%%%%%%%%%%%%%%%%%%%%%%

\begin{questions}

\addpoints

\question[10]\mbox{}
\textbf{TRUE or FALSE!}  Write  TRUE if the statement is true.  Otherwise, write FALSE.  Your response should be in ALL CAPS.  No justification is required.
\begin{enumerate}[(a)]
\item  The number $2$ is a prime element of the ring $\mathbb Z[\sqrt{-3}]$
\vspace{1.3in}
\item  There are exactly five Abelian groups of order $16$ up to isomorphism
\vspace{1.3in}
\item  In a ring $R$, if $ab=0$ then $a=0$ or $b=0$
\vspace{1.3in}
\item  There are two non-isomorphic non-Abelian groups of order $8$
\vspace{1.3in}
\item  The element $1+\sqrt{-3}$ is a unit in $\mathbb Z[\sqrt{-3}]$
\vspace{1.3in}
\end{enumerate}

\newpage
\question[10]\mbox{}
\begin{enumerate}[(a)]
\item  State the definition of an integral domain
\vspace{2in}
\item  State the Third Isomorphism Theorem for groups
\vspace{2in}
\item  State Cayley's Theorem
\end{enumerate}

\newpage
\question[10]\mbox{}

\begin{enumerate}[(a)]
\item Write down the definition of an irreducible element of an integral domain $R$
\vspace{2in}
\item Prove that $6+i$ is an irreducible element of the Gaussian integers $\mathbb Z[i]$
\end{enumerate}


\newpage
\question[10]\mbox{}

Let $N$ be a normal subgroup of $G$ with $[G:N] = n$.
\begin{enumerate}[(a)]
\item State Lagrange's Theorem
\vspace{2in}
\item Prove that for all $x\in G$, the order of $xN$ in $G/N$ divides $n$
\vspace{2in}
\item Prove that $x^n\in N$ for all $x\in G$.
\end{enumerate}

\newpage
\question[10]\mbox{}

The \textbf{ring of power series} $\mathbb R[[x]]$ is the set of all formal series expressions $\sum_{n=0}^\infty a_nx^n$, where $a_0,a_1,a_2,\dots$ are real numbers, with addition given by
$$\left(\sum_{n=0}^\infty a_nx^n\right) + \left(\sum_{n=0}^\infty b_nx^n\right) = \sum_{n=0}^\infty (a_n+b_n)x^n$$
and multiplication given by
$$\left(\sum_{n=0}^\infty a_nx^n\right) \left(\sum_{n=0}^\infty b_nx^n\right) = \sum_{n=0}^\infty c_nx^n,$$
where here $c_n = a_nb_0 + a_{n-1}b_1 + \dots + a_1b_{n-1} + a_0b_n$.
\begin{enumerate}[(a)]
\item Show that in this ring $1-x$ is a unit.  [Hint: geometric series $1+x+x^2+\dots=\frac{1}{1-x}$]
\vspace{2in}
\item Show that in this ring $x$ is not a unit.
\end{enumerate}


\newpage
\question[10]\mbox{}
Let $G$ be a finite group.
The \textbf{exponent} of $G$, denoted $\exp(G)$, is the maximum of the orders of elements of $G$.
\begin{enumerate}[(a)]
\item State the Structure Theorem for Finite Abelian Groups in Invariant Factor Form
\vspace{2in}
\item If $G$ is a finite Abelian group with invariant factors $a_1,a_2,\dots, a_r$ (in increasing order), prove that $a_r$ is the exponent of $G$
\vspace{4in}
\item Determine $\exp(G)$ for $G=\mathbb{Z}_{18}\times\mathbb{Z}_{21}$. [Hint: put it in invariant factor form]
\end{enumerate}

\newpage
\question[10]\mbox{}

Let $G$ be a group, suppose $g\in G$ is an element of order $n$, and let $H$ be the subgroup of $G$ generated by $g$.
Consider the group homomorphism
$$\varphi: \mathbb Z\rightarrow G,\quad f(k) = g^k.$$
\begin{enumerate}[(a)]
\item State the First Isomorphism Theorem for groups
\vspace{2in}
\item Determine the kernel of $\varphi$
\vspace{2in}
\item Explain why $\mathbb Z/\ker(\varphi)$ is isomorphic to $H$
\end{enumerate}



\end{questions}
\end{document}
